%% ============================================================================================
%% DRAFT v2, part 5: Threats to validity and Conclusion.
%% ============================================================================================

\section{Threats to validity}
\label{sec:threats}

We organise the threats by the four validity types used for case studies and experiments in software engineering
\citep{runeson2009guidelines,wohlin2012experimentation}; the threats specific to mining GitHub and to code-search
discovery, and how the design addresses them, are mapped in Table~\ref{tab:perils}.

\subsection{Construct validity}

\paragraph{Intended mode.} The construct behind the RQ2 contrast is the \emph{intended mode} of a delegate, annotated by
a language model from the \code{name} and \code{description} only; the body is not seen, so its construct validity
rests on the description stating the purpose. Agreement between the two model raters is below Krippendorff's threshold for
tentative conclusions on both the role and the mode (\S\ref{sec:annotation}), and a human
validation of \nHumanSample{} items has not yet yielded agreement estimates. A third model rater on the \nOpusItems{} validation items
(\S\ref{sec:annotation}) agrees with Sonnet more than with Haiku, on the mode (\alphaOpusModeSonnet{} against
\alphaOpusModeHaiku{}; Haiku and Sonnet \alphaKitModeHaikuSonnet{} on the same items) and on the role
(\alphaOpusRoleSonnet{} against \alphaOpusRoleHaiku{}), so the primary rater's mode labels are the noisiest of the
three. We draw no conclusion from the role labels. We could check the dependence of the mode contrasts on the rater
only partly. On the \nRaterSensSpecs{} engineered-population specifications in the reliability sample, the paired gap
is \RaterSensGapWriteHaiku{} for file-writing tools and \RaterSensGapAuthorityHaiku{} for any reach under Haiku's
labels (\RaterSensReposHaiku{} repositories hold both modes), and \RaterSensGapWriteSonnet{} and
\RaterSensGapAuthoritySonnet{} under Sonnet's (\RaterSensReposSonnet{} repositories). The two sets of repositories
differ and are small, and even under Haiku's labels the subsample does not reproduce the full-population gap, so the
check supports the direction of both gaps, and that the any-reach gap is much the smaller one, not their size. Misclassification that
is non-differential would attenuate both gaps, but the mode has four values and read-only wording in a description may
shape both the label and the tool list, so misclassification may be differential and we do not rely on attenuation.
The dominance measures do not use the mode labels.

\paragraph{What depends on model annotation.} The dominance measures are stable across the four populations,
including ALL, which uses no README classification (\S\ref{sec:res-rq3}). E2 depends on codebook P (one model rater; the third model agrees with it at
\alphaOpusEngHaiku{} on engineered against not, and among the \nEngKit{} sampled repositories calls \nEngOpusOnly{}
engineered that Haiku excluded and \nEngHaikuOnly{} the reverse, so the screen errs towards exclusion),
and the prose-restriction and shell-use results on codebooks S and B. The role labels mix two raters and two codebook
versions (\S\ref{sec:annotation}). Codebook K was applied
by Sonnet, which is also one of the models whose answers it judges; every positive and ambiguous label was re-read by the first author
against the answer (\S\ref{sec:res-rq5}).

\paragraph{Effective grant.} Capability is measured as the tool set the runtime resolves: our parser only extracts the
field's value, which is re-emitted verbatim into a probe file and resolved by the tool, so parser error is confined to
frontmatter extraction. The measure is still one of \emph{permission}; authority is argued analytically
(\S\ref{sec:capmodel}) and tested by execution (\S\ref{sec:res-rq5}). Whether it is realised depends on the
session's permission mode and user-level settings, which are unobservable. Deny rules and hooks in committed settings
can block matching commands; we did not evaluate which writes they block. Counting scoped shells as non-dominating
and omitting capability classes (egress, MCP tools, delegation, confidentiality) make our authority measures lower
bounds.

\paragraph{Prose restrictions and shell use.} Directive sentences are detected by a published lexicon that also counts
weaker negative markers as prohibitions (\S\ref{sec:nlp}); non-Latin bodies (\pctNonLatin{}) are excluded from that
stage, and restrictions stated without a deontic marker are missed, so topic shares describe explicit directives only.
Whether a specification is read-only is judged from its whole text instead (codebook S). Codebook B's agreement
between Haiku and Sonnet is below the threshold (\S\ref{sec:annotation}); the third model agrees with Haiku at
\alphaOpusShellHaiku{} (\alphaOpusShellChangeHaiku{} for change against the rest), and on codebook S at
\alphaOpusRestrHaiku{} for read-only against not. Codebook B's categories depend on what the text
mentions: a specification may need commands it does not name. Long bodies are truncated for annotation. The
annotators read untrusted repository text; we did not check it for instructions aimed at an annotator.

\paragraph{Engineered projects.} ``Engineered'' is a proxy for a construct with no agreed definition
(\S\ref{sec:curation}), which is why the grant results are reported under four nested populations.

\subsection{Internal validity}

\paragraph{Version conditionality.} Every resolution, admission and defect label is relative to one build (v2.1.233,
macOS arm64); \nRemovedNames{} names that earlier releases accepted no longer resolve, and the \code{Grep}/\code{Glob} behaviour is that of native builds from v2.1.160 onward; we did not re-run
the oracles on the later experiment builds.

\paragraph{Experiment.} The experiment establishes what a configuration permits in a controlled setting, not what
happens in production. It ran in \code{auto} mode, so the results for configurations with a shell are conditional on
commands running without per-command approval. Indirect requests through content the agent reads were not tested. The tasks are
simple and benign and the fixtures synthetic. The system prompt includes the harness note, which asks for a concise,
literal final answer and may affect how answers are phrased, including the completion claims. The runs span three
builds and two orchestrating sessions whose context is not part of the stimulus but is not independent across runs
within a session; we detect no difference between builds, but with \ExpBuildCellSizes{} runs per cell and build only
large differences were detectable (\S\ref{sec:res-rq5}). The scoped-shell probe ran in \code{auto} mode only; whether a
command pattern in the \code{tools} field acts as a pre-approval that leaves other commands to prompt in
\code{default} mode was not tested.

\paragraph{Reconstruction.} \pctUnavailable{} of inventoried files could not be fetched; all but
\nFilesObjectMissing{} belong to repositories that are no longer available, so the loss is almost entirely of whole
repositories rather than selected files.

\subsection{External validity}

The frame is about \pctFrameOfRerun{} of the repositories that code search could find, selected by an undocumented
ranking (\S\ref{sec:collection}), so the results describe the frame. File-level discovery biases team size upward;
the dominated-withholding share shows no trend across repository sizes (\S\ref{sec:res-rq3}), but we cannot rule out
other selection effects. The corpus is Claude Code's
format at one snapshot. The mechanism generalises in principle, since Copilot and Codex expose comparable subagent files
and Cursor reads the same directory \citep{GitHubCustomAgentsDoc,OpenAICodexSubagentsDoc,CursorSubagentsDoc}; the rates
do not, and private use is invisible to us.

\subsection{Conclusion validity}

Specifications within a repository are dependent (intra-class correlation \iccValue{} for whether an explicit grant
can reach the file system, design effect \deffValue{}), so
prevalences use the repository as the unit, with a bootstrap over repositories, and role contrasts are paired within
repository; descriptive shares of corpus composition are pooled and labelled as such. The prompt contrasts are
pre-specified, Holm-adjusted \citep{holm1979simple} and read together with the per-model rates. Annotation and agents
are non-deterministic, so annotation is repeated by a second model on reliability samples and each experiment cell
\nRepsPerCell{} times \citep{ouyang2025nondeterminism,baltes2026guidelines}.

\paragraph{Duplication.} Byte-identical specifications recur across repositories (\S\ref{sec:res-rq1}), so
repositories are not independent observations. Keeping one copy per content hash (\nSpecsDedup{} of \nSpecsEng{}
specifications) leaves the headline results unchanged: \DedupImplicit{} omit \code{tools}, \DedupIncoherent{} of
write-withholding explicit grants keep an unscoped shell, and the paired gap is \DedupGapWrite{} for file-writing tools
and \DedupGapAuthority{} for any reach.

\section{Conclusion}
\label{sec:conclusion}

This paper contributes a snapshot-exact corpus of \nSpecsAll{} subagent specifications, executable oracles that obtain
what the coding tool loads and grants, model annotations of role, intent and prose restriction with published
codebooks, a controlled experiment of \nExpRuns{} delegations with a report-only extension, and a linter for the
defects it measures.

In \nReposEng{} engineered repositories developers define \nSpecsEng{} delegates, and they restrict them: agents meant
to inspect receive file-writing tools far less often than agents meant to change things. But in the average
repository \RQthreeIncoherentEng{} of explicit grants that withhold every file-writing tool keep an unscoped shell, and
the text of most of these specifications asks the shell for at most reading or verification (\ShellUseNoChangeNeed{}
never ask for a change by command). Such grants withhold permission, not authority, wherever shell commands run
unapproved and unsandboxed.

Repeated execution separates what a grant permits from what agents do. In \code{auto} mode, with the most common
inspection grant and no prompt restriction, requested changes succeeded in \ExpCtwoRate{} of runs, and without a shell
in \ExpCfourRate{}. In the realistic inspection configuration a prompt restriction held for Sonnet and Opus but not
reliably for Haiku (\ExpCthreeRange{} across models). On one fixture with report-only tasks, no agent changed a file.
Within these limits, the risk we observed came from changes that were requested and not refused, not from agents
acting unasked.

Developers also get no feedback on these files: \RQsixUnresolvedRepos{} of the repositories that list tools name a tool the runtime
does not resolve, most files naming a removed tool were committed after its removal, and duplicate names discard
\RQsixShadowed{} specifications without a message. Future work should test whether command-scoped shells in a subagent's \code{tools}
field restrict anything outside auto mode, how the results change under other permission modes and under prompt injection, whether the pattern
holds in other harnesses, and whether warnings change what developers write.
